\chapter{Quantum Mechanics}
\label{chap:quantum-mechanics}

\section*{Introduction}

Quantum mechanics is extraordinarily successful, and one of the most tested theories in all of science.
In its modern, field-theoretic form, it does not fundamentally describe particles, but rather the evolution of a mathematical object called the \emph{wavefunction} (or, in relativistic theory, a quantum field). 
What we perceive as particles emerge as localized outcomes of interactions and measurements.
(This is an interpretive stance shared by most, but not all, approaches; Bohmian mechanics, for example, retains particles with definite positions.)
Once the amplitude and phase of a system's wavefunction are known, the wavefunction's future can be predicted with astonishing precision. 
The outcomes of measurements, however, can be predicted only probabilistically.
Particle properties such as mass, charge, and spin are encoded in the wavefunction, and its dynamics are governed by the Schrödinger equation.


\section*{Particles}

In experiments, we observe discrete events---localized detections that we call particles. Elementary particles come in two broad classes, distinguished by their spin:

\begin{itemize}
\item \textbf{Fermions}: half-integer spin; the constituents of ordinary matter
\item \textbf{Bosons}: integer spin; this class includes the force carriers
\end{itemize}

Fermions include quarks and leptons. There are six types of quarks and six leptons (electron, muon, tau, and their corresponding neutrinos). These form the building blocks of matter.

Bosons include the photon, gluons, and the $W$ and $Z$ bosons, which mediate the fundamental interactions, as well as the Higgs boson, associated with the field that gives mass to the $W$, $Z$, and the elementary fermions. Note that the boson/fermion distinction is about spin and statistics, not about ``matter versus force'': composite particles such as mesons or helium-4 atoms are bosons, yet they are made of matter. Note also that being ``virtual'' is not a property of a class of particles; virtual particles are a calculational device describing internal lines in perturbation theory, and they can be bosons or fermions.

Despite their diversity, elementary particles are characterized by only a few properties: mass, charge, and spin.

Interestingly, not all of these properties are ``quantized'' in the same sense.

Spin takes only integer or half-integer multiples of $\hbar$, a fact tied directly to the symmetry of rotations. Electric charge is also observed to come in discrete units: free particles carry integer multiples of the electron charge $e$, while quarks carry $\pm e/3$ or $\pm 2e/3$ but are never observed in isolation. Unlike spin, however, the quantization of charge is not explained by the Standard Model; it remains an open question, with Dirac's magnetic monopole argument and grand unified theories among the proposed explanations. Mass is different again. While each elementary particle has a well-defined mass, there is no known fundamental unit from which all masses are built. The masses of the elementary particles arise from their interaction with the Higgs field, and the pattern of their values remains one of the open questions in fundamental physics. (Most of the mass of ordinary matter, such as protons and neutrons, does not come from the Higgs mechanism at all, but from the binding energy of the strong interaction.) So quantization is not a single universal mechanism, but can arise in different ways from the underlying structure of the theory.

The spin of fermions has an interesting property. It describes how a particle's state transforms under rotations. A striking consequence is that the state vector of a fermion
(spin $1/2$) acquires a minus sign after a $360^\circ$ rotation and returns to itself only after a full $720^\circ$ rotation, while that of a boson (integer spin) returns after $360^\circ$. The sign is a global phase and is not directly observable for an isolated particle, but it becomes physical in interference experiments against a reference.

From this small set of elementary particles, the entire visible universe is constructed. The richness of physical reality emerges from surprisingly minimal ingredients. 

\subsection*{Emergence}

Composite systems can exhibit properties and behavior not present in their individual constituents. 
For example, an atom containing an even number of fermions (protons, neutrons, and electrons together) behaves as a boson, while one containing an odd number behaves as a fermion. 
This determines the quantum statistics of collections of such atoms: whether they obey the Pauli exclusion
principle or can occupy the same quantum state.

The same principle extends to higher levels of organization. Atoms, molecules, and macroscopic objects exhibit collective properties that
cannot be understood simply by examining their constituents in isolation. New behavior and effective rules emerge through their
organization and interactions.

Classical mechanics provides another example of emergence. By Ehrenfest's theorem, the expectation values of quantum position and momentum evolve according to 
equations resembling Newton's equations. Under appropriate conditions (for example, when the force varies slowly across the extent of the wavepacket), classical particle dynamics emerges as
an approximation to the underlying quantum description.


\section*{The Wavefunction}

The wavefunction is the central object of quantum mechanics. It encodes all physically accessible information about a system.

Mathematically, the state of a system is represented by a vector
\[
|\psi\rangle \in \mathcal{H},
\]
where $\mathcal{H}$ is a Hilbert space.

Its time evolution is governed by the Schrödinger equation:
\[
i\hbar \frac{\partial}{\partial t} |\psi(t)\rangle = \hat{H} |\psi(t)\rangle.
\]

This evolution is linear and unitary. At this level, the evolution of the wavefunction is fully deterministic.

Physical quantities are represented by operators acting on the state, and measurement outcomes correspond to their eigenvalues.

\subsection*{Measurement}

We do not observe the wavefunction directly. Instead, we observe specific results---localized events in spacetime.
The probabilities of these outcomes are given by the Born rule, a postulate of the theory:
\[
P(x,t) = |\psi(x,t)|^2 .
\]
Here $P(x,t)$ is a probability \emph{density}; the probability of finding the particle in a region is obtained by integrating it over that region.

The values of the wavefunction are complex numbers, and complex numbers cannot be interpreted as probabilities in any meaningful way.
The Born rule acts as a translator that turns them into probabilities by taking the absolute square \( |\psi|^2 \).
The result is a non-negative, real number that specifies the probability density. 

The double-slit experiment illustrates this.
Even when particles are sent one at a time, an interference pattern emerges.
The wavefunction propagates through both slits, while detection produces localized impacts.

This leads to a dual structure:

\begin{itemize}
\item Continuous, linear, unitary evolution of the wavefunction
\item Discrete, localized outcomes upon measurement, with probabilistic predictions
\end{itemize}


\subsection*{Superposition as Informational Encoding}

A quantum system can exist in a superposition of states:

\[
|\psi\rangle = \alpha |\text{H}\rangle + \beta |\text{T}\rangle,
\quad
|\alpha|^2 + |\beta|^2 = 1.
\]

This does not mean the system is partly in each state in a classical sense. Rather, the wavefunction encodes multiple possibilities simultaneously in a single mathematical object.

To make this intuitive, consider a quantum coin. Unlike a real coin in our geometric space, a quantum coin exists as a superposition of states
within a two-dimensional Hilbert space. 
The coefficients $\alpha$ and $\beta$ are \textbf{probability amplitudes},
represented by complex numbers. They are not probabilities themselves;
rather, the probability of each outcome is determined by the absolute
square (the Born rule):

\[
P(H)=|\alpha|^2,\qquad P(T)=|\beta|^2.
\]

The state must be normalized:

\[
|\alpha|^2 + |\beta|^2 = 1.
\]

For a fair quantum coin, the two outcomes are equiprobable, so the
amplitudes have equal magnitudes:

\[
|\alpha|=|\beta|=\frac{1}{\sqrt{2}}.
\]

The phases of the amplitudes, which do not affect these probabilities, become important when there are \emph{several indistinguishable ways} to reach the \emph{same} outcome. In that case, amplitudes are added before probabilities are extracted. If a particle can reach a detector via two paths with amplitudes $\psi_1$ and $\psi_2$, then
\[
P = |\psi_1 + \psi_2|^2 = |\psi_1|^2 + |\psi_2|^2 + 2\,\mathrm{Re}(\psi_1^*\psi_2).
\]
The last term, which depends on the relative phase, is the interference term. This phase-sensitive structure is what produces the fringes in the double-slit experiment.


\subsection*{Missing Identity}

Particles do not have identity, that is, a property that would allow us to follow them as distinct individuals.  
In classical physics, even two perfectly identical objects can be told apart by tracking their trajectories: this is the one that went left, that is the one that went right.
Quantum mechanics removes the trajectory, and with it, the basis for identity.

Consider a parabolic mirror. A photon emitted from the focal point is later detected on a screen after reflecting from the mirror. Classically, 
one would imagine the photon taking a specific path: leaving the source, reflecting at a definite point, and arriving at the detector.
However, the photon is not assigned a single trajectory. Instead, it can be shown that all the possible paths it can take contribute to the outcome. 
Contributions from the entire surface of the mirror determine where the photon most likely hits the screen. The emission and detection events are well defined, 
but between them there is no single path to follow, only a wavefunction giving the probability of where the photon will be detected. 
Without a path, there is nothing that could persist as ``this particular photon'' between the source and the detector.

The same holds for tunneling. A particle that is incident on a barrier it classically could not cross can nevertheless be detected on the other side, because the wavefunction decays within the barrier but does not vanish. 
There is no trajectory that carries the particle ``through'' the barrier, and so no meaningful answer to the question of where it was, or how fast it traveled, while inside. 
(Apparently superluminal tunneling times, known as the Hartman effect, refer to a delay of the wave envelope; they cannot be used to send signals faster than light.)

This has a precise mathematical consequence. If identical particles cannot be tracked, then exchanging them cannot produce a physically different state: the two states can differ at most by a phase, which can only be $+1$ or $-1$. This is exactly the exchange symmetry described in the next section. 
Strictly speaking, particles do not ``not exist''; detection events are real, and in quantum field theory particles are excitations of a field. What is missing is their persistent, individual identity. 
A better picture is that of a field, with quanta that are counted but not labeled.


\subsection*{Pauli Exclusion Principle}

In Hilbert space, exchanging two identical particles corresponds to a transformation of the state:

For bosons:
\[
\psi(x_1, x_2) = +\psi(x_2, x_1)
\]

For fermions:
\[
\psi(x_1, x_2) = -\psi(x_2, x_1)
\]

Here $x$ stands for all the degrees of freedom, including spin. For two fermions in the same spin state, the antisymmetry implies that the spatial wavefunction satisfies:
\[
\psi(x,x) = 0 .
\]

This is the Pauli exclusion principle: no two identical fermions can occupy the same quantum state. (Two electrons with opposite spins can share the same spatial orbital, because the antisymmetry is then carried by the spin part of the state.)

For bosons, symmetric states allow multiple occupation. Any number of photons, for example, can occupy the same mode of the electromagnetic field. 
This enables phenomena such as Bose--Einstein condensation, where a macroscopic number of bosonic atoms (for example rubidium-87) occupy the same quantum state, leading to collective quantum behavior.


\subsection*{Heisenberg Uncertainty}

A plane wave:
\[
\psi(x) = e^{ikx}
\]
has a well-defined momentum $p=\hbar k$ but is completely delocalized. (Strictly, such a wave is not normalizable and is an idealization.)

To localize a particle, one must superpose many momenta:
\[
\psi(x) = \int a(k) e^{ikx} \, dk
\]

The more localized the position, the broader the momentum distribution. This is expressed by the uncertainty principle:
\[
\Delta x \Delta p \ge \frac{\hbar}{2}.
\]

It would be easy to think that the uncertainty is just the problem of observation and that particles would actually have both well defined position and momentum, simultaneously. 
However, this is not the case.

Consider a single free particle in a vast space. Its wavefunction is a pure sine wave. Its energy is related to its frequency by $E=\hbar\omega = hf$. Its probability density is perfectly flat. 
Because space is so huge, the probability of finding the particle in any one small region is practically zero. In this state, the particle's position is totally smeared out---it simply doesn't 
have a specific `here' or `there'. 
The only way to localize the particle is to stack additional waves into the wavefunction. 
These waves interfere with each other, canceling out in most places and peaking in one small spot. 
But then you are mixing multiple momenta, and therefore multiple energies, together. The more you pin down the position, the more the momentum `smears.' The fuzziness isn't a measurement error; it's just how waves work.

A familiar analogy is a musical note, which illustrates the closely related trade-off between time and frequency: you can have a pure, steady pitch that lasts forever (perfect frequency, no timing), or you can have a short, sharp `staccato' pop (perfect timing, no pitch). Since $E=\hbar\omega$, this is an energy--time trade-off, $\Delta E\,\Delta t \gtrsim \hbar/2$. (It differs in character from $\Delta x\,\Delta p$, because time is a parameter in quantum mechanics and not an operator.)
Either way, it is a property of waves in general that you cannot have both at once.


\subsection*{Entanglement}

Composite systems can exhibit correlations that cannot be reduced to their parts.
Such correlations violate Bell-type inequalities, which rules out local hidden-variable explanations. 
(Entanglement does not, however, allow signaling faster than light.)

For example:
\[
|\Psi\rangle = \frac{1}{\sqrt{2}}
\left(
| \uparrow \downarrow \rangle - | \downarrow \uparrow \rangle
\right)
\]

This state cannot be written as a product of individual states. The system must be described as a whole.

Entanglement is a defining feature of quantum mechanics, linking subsystems into a single, inseparable structure.


\section*{Quantum Field Theory (QFT)}

When quantum mechanics is combined with special relativity, the natural framework becomes quantum field theory.

In QFT, fields---not particles---are fundamental. Every point in space is associated with field values. Particles appear as excitations of these fields.

A useful analogy is a continuous medium composed of coupled oscillators. Disturbances propagate as waves, and quantized excitations are observed as particles.

For a scalar field $\phi(x)$ with conjugate momentum $\Pi$, the energy density may take the form:
\[
\mathcal{H} = \frac{1}{2} \Pi^2 + \frac{1}{2} (\nabla \phi)^2 + \frac{1}{2} m^2 \phi^2
\]

Different fields correspond to different particles. Interactions arise from couplings between fields.


\section*{The Measurement Problem}

Quantum mechanics suffers from the so-called measurement problem. 
QM describes a superposition of possibilities, implying that a system is in a combination of all possible states, such as here and there, simultaneously, until observed.

But what does it mean to observe? What constitutes an ``observer''?

Different interpretations offer different answers:

\begin{itemize}
\item Copenhagen: measurement introduces classical outcomes via so-called ``collapse of the wavefunction''
\item Many-Worlds: all outcomes occur in branching states
\item Objective collapse: wavefunction collapse is a physical process
\item Relational / information-based views (e.g.\ relational quantum mechanics, QBism): states are observer-dependent or reflect an agent's information
\end{itemize}

A modern operational view is that an observer is any system that produces a stable record of information through interaction.

% Need to hypothesize whether my wife meets the criteria of being a stable record of information.
% That is a brave hypothesis to test! Proceed with extreme caution---that particular variable has a
% reputation for being highly volatile, regardless of what the laws of physics might suggest!


\subsection*{Decoherence and Irreversible Recording}

When a system interacts with an environment, entanglement spreads information across many degrees of freedom.

At the level of the full system, evolution remains unitary. However, when one considers only a subsystem, information appears to be lost.

This is captured by the reduced density matrix:
\[
\rho_S = \mathrm{Tr}_E(\rho_{SE})
\]

Interactions with the environment suppress interference terms---a process known as decoherence.

The entropy of the subsystem typically increases:
\[
S(\rho) = -\mathrm{Tr}(\rho \log \rho)
\]

Information becomes distributed into many inaccessible degrees of freedom. Reversing this process would require control over an enormous number of variables.

Thus, irreversibility emerges not as a fundamental law, but as a practical consequence of complexity.

An ``observer,'' in this sense, is a system that participates in such irreversible encoding of information.
Decoherence explains why interference becomes unobservable and why classical-looking records form, but by itself it does not explain why a single outcome is obtained; that is the part on which the interpretations above differ.


\subsection*{What Is the Quantum System?}

Everything described above is an observed fact about the universe in which we live.

Consider a real-world object---say, a car engine. We observe a system composed of pistons, gears, a carburetor, and so on. 
Sooner or later, by examining its structure and behavior, we are able to figure out its purpose.

With quantum mechanics we are in a similar position, but with the mathematics only. 
We have observed a system described by a wavefunction, and we have cataloged its properties. 
We have observed that this wavefunction evolves according to the Schrödinger equation, and we have observed that outcomes appear as if it collapses upon measurement.

What exactly is the system that all of this describes? 

What is it that we are actually looking at?